\documentclass[11pt]{article}
    \usepackage[utf8]{inputenc}
    \usepackage[T1]{fontenc}
    \usepackage[margin=1in]{geometry}
    \usepackage{hyperref}
    \usepackage{enumitem}
    \title{Teaching Indi | to Turn}
    \author{Piter Garcia}
    \date{March 20, 2026}
    \begin{document}
    \maketitle
    \section*{Quick Scan}
    \begin{itemize}[leftmargin=*]
    \item Grade: Kindergarten
    \item Workbook sessions: Session 24
    \item Class blocks: Day 1 12:25-1:15 Justinger, Day 2 12:25-1:15 Kesys, Day 4 12:25-1:15 Pum
    \item Reference file: teaching-indi-to-turn.bib
    \end{itemize}
    \section*{School Planning Goals and Inclusion}
    \begin{itemize}[leftmargin=*]
    \item What is expected: I can identify key parts of Indi and what they do (K-1.NSD.2) | I can teach someone else about Indi and how it works (using color mats for input & output) (K-1.NSD.1) | I can show Indi how to travel in different directions (K-1.NSD.1) | I can follow an algorithm to complete a task (K-1.CT.6) | I can makes changes to my program when it is not working (K-1.CT.9) | I can program Indi to travel along different paths and around objects (K-1.CT.10)
    \item What we achieved: No direct transcript match is attached yet. Treat this lesson as workbook-backed and enrich it when a matching classroom transcript is available.
    \item What needs improvement: stronger proof, cleaner assessment notes, and tighter lesson artifacts.
    \item How to improve: add transcript proof, one saved artifact, and clearer success criteria.
    \item Inclusion focus: use visual modeling, chunked directions, predictable routines, and flexible participation roles.
    \end{itemize}
    \section*{Official References}
    \begin{itemize}[leftmargin=*]
    \item \url{https://csteachers.org/k12standards/interactive/}
\item \url{https://dmmedia.sphero.com/email-marketing/Sphero-Edu/SpheroEdu-k12-teacher-resource-guide-v1_updated050818.pdf}

    \end{itemize}
    \end{document}
